\section{Certified Installation Details}
\label{app:transfer-details}

Execution signers durably retain canonically encoded, materialized deltas and transaction
outputs before signing and serve them by digest. The delta includes writes,
deletions, statuses, events, fees, and authenticated consumed-ID and processing metadata, not private worker IDs
or physical database layout. A receiver can fetch one complete authenticated copy
from a signer or another peer; it need not download \(f+1\) copies.

The receiver checks full ordering finality, signer eligibility, commitments, and
the exact parent history and processing boundary. It applies verified data to
staging state and checks the post-state root before crash-safe atomic publication
of state, outputs, consumed-ID history, and progress. A checkpoint must preserve
equivalent replay protection before earlier deduplication history is discarded.
Duplicate installation is a no-op. Partial transfers and stale
worker completions must not overwrite published state; caches for later cuts are
revalidated against the installed parent. These are required properties, not an
implemented crash-recovery protocol.

If the parent is unavailable, the receiver must recover a certified chain of
updates or an authenticated checkpoint; equal roots at different processing
positions are insufficient. Retention, checkpoint availability, garbage
collection, and independent endorsement discovery remain obligation D2 in the
main paper. Recovery requires eventual synchrony, fair service, and sufficient
capacity. When certified updates cannot be retrieved, re-executing authenticated
inputs from a verified parent is a possible fallback, not a guarantee of bounded
catch-up time.
